\subsection{结论}\label{sec:19d}
我们将在 Fukaya 范畴框架内进行类似的代数构造，以取代从 Picard--Lefschetz 数据几何地恢复各 $X^k$ 的尝试。事实表明，在这一层次上，上述困难消失了，比较\stref{笔记}{16.11}。为简化叙述，将 $\Psi^k$ 与 $\varpi^k$ 分别当作真正的正合 Lefschetz 纤维化与双纤维化。这样做没有问题：$\Psi^k$ 的情形已在\secref{19b}解释，$\varpi^k$ 的讨论完全类似。因此，有两列 Fukaya 范畴 $\F(X^k)$ 与 $\F(\Psi^k)$。对每个 $0<k\le n$，令 $V_j^k\subset X^k$ 为 $\gamma_j^{k-1}$ 对应的消失圈；由构造，它也是 $\mu_j^k$ 对应的匹配圈，见\stref{引理}{16.15}。照常将它们赋予 Lagrangian 膜结构，将所得对象集记为 $V^{k,\#}$。写成
\[
\A^k=\F(V^{k,\#}),\qquad
\A^{k,\to}=\F^\to(V^{k,\#});
\]
也就是说，$\A^k$ 是 $\F(X^k)$ 中由这组消失圈构成的全 $\Ai$ 子范畴，$\A^{k,\to}$ 为相应有向 $\Ai$ 子范畴。

算法的起点与前面相同，即在拓扑上重构 $X^n$、$\mathcal K_{X^n}$ 的平凡化及消失圈 $V^n$。给定这些数据后，可按第\sr{section:13}{13 节}的方法组合地计算 $\A^{n,\to}$。迭代步骤如下。假设给定一个拟同构于 $\A^{k+1,\to}$ 的有向 $\Ai$ 范畴 $\B^{k+1,\to}$。由\stref{命题}{18.13}与\stref{定理}{18.24}，得到在上同调上全忠实的 $\Ai$ 函子
\begin{equation}\label{eq:19.6}
\F(X^k)\longrightarrow\F(\Psi^k)\longrightarrow\Tw\A^{k+1,\to}.
\end{equation}
如\stref{推论}{18.26}的证明中那样，可确定任意匹配圈在\eref{19.6}下的像，以及 $\Tw\B^{k+1,\to}$ 中的相应对象，可能相差一个分次平移；这种不确定性最终无关紧要，为简便，假设分次已正确选取。因此，$\A^k$ 拟同构于一个显式给出的全 $\Ai$ 子范畴 $\B^k\subset\Tw\B^{k+1,\to}$。若 $k>1$，就只保留拟同构于 $\A^{k,\to}$ 的 $\B^{k,\to}$，\emph{舍去部分信息}，再继续到高一维。这样到达 $k=1$ 时，得到拟同构于 $\A^1$ 的 $\Ai$ 范畴 $\B^1$。最后还可再前进半步，重新应用\eref{19.6}，说明 $\F(X)$ 本身拟等价于 $\Tw\B^{1,\to}$ 的某个事先未知的全 $\Ai$ 子范畴。\znote{302}{为什么要舍去：归纳的输入只有有向范畴 $\A^{k,\to}$，而 $\B^k$ 中非有向的态射与高阶运算在\eref{19.6}的函子里根本不出现；保留它们反而使下一步无法与 $\Tw\B^{k+1,\to}$ 比较。这也是上面 19c 末尾“数据不足以恢复 $X^k$”的代数对应物：几何上丢掉的是同痕，代数上丢掉的是非有向部分。[来源：A，原书本节与 19c 末段。]}

\SourcePage{311}
\begin{proof}[定理~\sr{st:B}{B} 在 $n\ge1$ 时的证明]
由上述论证，结论现已显然成立，所考虑的 $\Ai$ 范畴为 $\Tw\B^{1,\to}$。
\end{proof}

\begin{proof}[推论~\sr{st:C}{C} 在 $n\ge1$ 时的证明]
由\stref{推论}{5.27}，可将 $\B^{1,\to}$ 替换为一个与之拟同构、实际上是 $\dg$ 范畴的有向 $\Ai$ 范畴 $\C^\to$。取在上同调上全忠实的 $\Ai$ 函子 $\F(X)\to\Tw(\C^\to)$。必要时引入已有对象的重复副本后，可假设此函子在对象上单射，于是 $\F(X)$ 拟同构于其像，而该像是具有有限维态射空间的 $\dg$ 范畴。
\end{proof}

\begin{proof}[定理~\sr{st:A}{A} 的证明]
由\stref{命题}{19.7}，消失圈分裂生成 $\Pi\Tw\F(X^1)$。由\stref{推论}{4.9}与\stref{引理}{4.7}，$\Pi\Tw\F(X^1)$ 拟等价于 $\Pi\Tw\A^1$，因此也拟等价于 $\Pi\Tw\B^1$。
\end{proof}

细心的读者会注意到，该算法在最低维时未给出完整信息：可以很容易地确定 $\F(X^n)$ 的任意有向 $\Ai$ 子范畴，但这一做法并不先验地适用于 $\F(X^n)$ 本身。部分原因在于，没有从真正平凡的零维情形开始归纳；可如下弥补。取 $\Psi^n:X^n\to\CC$，按第\sr{section:18}{18 节}的过程定义 $\F(\Psi^n)$，即取沿光滑纤维分歧的二重覆盖 $\widetilde X^n\to X^n$，并考虑 $\Z/2$-不变 Lagrangian 子流形。特别地，$\widetilde X^n$ 中有 (B) 型等变膜，它们是消失路径 $\gamma_j^n$ 相应 Lefschetz 顶针的原像。考察\stref{命题}{18.14}的证明，可知它们在 $\F(\Psi^n)$ 中构成例外集。因此，相应全 $\Ai$ 子范畴与有向 $\Ai$ 子范畴拟同构。另一方面，该有向 $\Ai$ 子范畴可用第\sr{section:13}{13 节}的组合方法处理，因为定义相关 Floer 上链群与 $\mu^d$ 运算时，可以使用标准复结构及零非齐次项，\stref{假设}{14.5}仍会满足。在这一特殊情形中，所有伪全纯多边形实际上都是常值的；其正则性来自\stref{引理}{11.5}，如\stref{引理}{14.6}、\stref{引理}{14.7}的证明中那样。将所得有向 $\Ai$ 范畴记为 $\B^{n+1,\to}$。它带有在上同调上全忠实的函子
\begin{equation}\label{eq:19.7}
\F(X^n)\longrightarrow\Tw\B^{n+1,\to}.
\end{equation}
利用它，可将定理~\sr{st:B}{B} 与推论~\sr{st:C}{C} 的证明推广到 $n=0$ 的情形，此时 $X^n=X$。

\StatementNumber{10}\begin{note}\label{st:19.10}
为类似地推广定理~\sr{st:A}{A}，需确定 $\Psi^{n-1}$ 的消失圈在\eref{19.7}下的像。遗憾的是，这里有一个高维时未出现的微小但棘手之处：在\stref{命题}{18.21}的对应陈述中，零次态射空间现在为二维，
\SourcePage{312}
必须在该空间的正确元素上构造锥，它对应于消失圈上的非平凡 Pin 结构。
\end{note}
